\documentclass[11pt, a4paper]{article}

% Gemeinsame Style-Definitionen (TEXINPUTS muss templates/ enthalten — siehe scripts/build_tex.py)
% bfw-style.tex — gemeinsame Style-Definitionen für alle Handouts/Aufgabenhefte
% Voraussetzung: Kompilation mit xelatex (wegen fontspec)

\usepackage[ngerman]{babel}
% Babel-ngerman macht " zu einem aktiven Shorthand (z.B. für "a -> ä). In unseren
% Texten verwenden wir " als normales Zeichen — daher Shorthand komplett ausschalten.
\shorthandoff{"}
\usepackage{geometry}
\geometry{a4paper, top=2.4cm, bottom=2.4cm, left=2.3cm, right=2.3cm,
          headheight=15pt, headsep=0.7cm, footskip=1.1cm}

\usepackage{fontspec}
% Fallback-Kette: Source Sans Pro -> Calibri -> Segoe UI -> Latin Modern Sans
% (Latin Modern ist immer mit MiKTeX vorhanden)
\IfFontExistsTF{Source Sans Pro}{%
  \setmainfont{Source Sans Pro}[Ligatures=TeX]%
}{\IfFontExistsTF{Calibri}{%
  \setmainfont{Calibri}[Ligatures=TeX]%
}{\IfFontExistsTF{Segoe UI}{%
  \setmainfont{Segoe UI}[Ligatures=TeX]%
}{%
  \setmainfont{Latin Modern Sans}[Ligatures=TeX]%
}}}
\IfFontExistsTF{Source Code Pro}{%
  \setmonofont{Source Code Pro}[Scale=0.92]%
}{\IfFontExistsTF{Consolas}{%
  \setmonofont{Consolas}[Scale=0.92]%
}{%
  \setmonofont{Latin Modern Mono}[Scale=0.92]%
}}

% Gesperrte Versalien für Labels/Titelbalken (Kapitälchen-Ersatz, funktioniert
% mit jeder OpenType-Schrift — Latin Modern Sans hat keine echten Kapitälchen)
\newcommand{\bfwlabelfont}{\footnotesize\bfseries\addfontfeatures{LetterSpace=8.0}}

\usepackage{xcolor}
% Farbpalette: hell, freundlich, modern
\definecolor{bfwPrimary}{HTML}{0EA5A4}    % Petrol/Türkis - Primär
\definecolor{bfwPrimaryDark}{HTML}{0F766E} % dunkler Petrol für Headings
\definecolor{bfwAccent}{HTML}{F59E0B}     % Amber - Akzent
\definecolor{bfwSuccess}{HTML}{10B981}    % Grün - Erfolg / Lösung
\definecolor{bfwWarn}{HTML}{F97316}       % Orange - Achtung
\definecolor{bfwInk}{HTML}{0F172A}        % fast Schwarz
\definecolor{bfwInkSoft}{HTML}{475569}    % gedämpftes Grau
\definecolor{bfwBgSoft}{HTML}{F8FAFC}     % off-white
\definecolor{bfwBgTeal}{HTML}{ECFEFF}     % helles Türkis
\definecolor{bfwBgAmber}{HTML}{FEF3C7}    % helles Gelb
\definecolor{bfwBgRose}{HTML}{FFE4E6}     % helles Rosé für Eselsbrücken
\definecolor{bfwTabHead}{HTML}{E4F3F2}    % zarte Kopfzeilen-Hinterlegung für Tabellen

\usepackage[colorlinks=true, linkcolor=bfwPrimaryDark, urlcolor=bfwPrimaryDark]{hyperref}
\usepackage{lastpage}
\usepackage{enumitem}
\setlist[itemize]{leftmargin=*, itemsep=2pt, topsep=2pt}
\setlist[enumerate]{leftmargin=*, itemsep=2pt, topsep=2pt}

\usepackage[most]{tcolorbox}
\usepackage{booktabs}
\usepackage{colortbl}
\usepackage{tikz}
\usetikzlibrary{decorations.pathmorphing, positioning, shapes.geometric, arrows.meta}

% ---------------------------------------------------------------------------
% Einheitliches Tabellen-Design
%   - etwas mehr Zeilenluft, dezente graue Linien (wirkt auch mit \hline ruhiger)
%   - Kopfzeile einer Tabelle bei Bedarf zart hinterlegen: \rowcolor{bfwTabHead}
% ---------------------------------------------------------------------------
\renewcommand{\arraystretch}{1.25}
\arrayrulecolor{bfwInkSoft!50}
\setlength{\heavyrulewidth}{1pt}
\setlength{\lightrulewidth}{0.5pt}

% ---------------------------------------------------------------------------
% Boxen — klare farbige Titelbalken mit gesperrten Versal-Labels
% (Titel bleibt per Option überschreibbar: \begin{merksatz}[title={...}])
% ---------------------------------------------------------------------------
\tcbset{bfwbox/.style={%
  enhanced, breakable,
  boxrule=0.5pt, arc=2.5pt,
  left=10pt, right=10pt, top=8pt, bottom=8pt,
  toptitle=4pt, bottomtitle=4pt,
  titlerule=0pt,
  fonttitle=\bfwlabelfont,
  coltitle=white
}}

% Merkkästchen — wichtige Definition / Kernpunkt
\newtcolorbox{merksatz}[1][]{%
  bfwbox,
  colback=bfwBgTeal,
  colframe=bfwPrimaryDark,
  colbacktitle=bfwPrimaryDark,
  title={MERKE},
  #1
}

% Eselsbrücke — Mnemoniks
\newtcolorbox{eselsbruecke}[1][]{%
  bfwbox,
  colback=bfwBgRose,
  colframe=bfwAccent!85!black,
  colbacktitle=bfwAccent!85!black,
  title={ESELSBRÜCKE},
  #1
}

% Beispiel-Box
\newtcolorbox{beispiel}[1][]{%
  bfwbox,
  colback=bfwBgSoft,
  colframe=bfwInkSoft,
  colbacktitle=bfwInkSoft,
  title={BEISPIEL},
  #1
}

% Lösungs-Box (für Aufgabenhefte)
\newtcolorbox{loesung}[1][]{%
  bfwbox,
  colback=bfwSuccess!7,
  colframe=bfwSuccess!65!black,
  colbacktitle=bfwSuccess!65!black,
  title={LÖSUNG},
  #1
}

% Achtung-Box — typische Fehler / Stolperfallen
\newtcolorbox{achtung}[1][]{%
  bfwbox,
  colback=bfwWarn!6,
  colframe=bfwWarn!85!black,
  colbacktitle=bfwWarn!85!black,
  title={ACHTUNG},
  #1
}

% Formel-Box — für zentrale Formeln (groß, ruhig, ohne Titelbalken)
\newtcolorbox{formelbox}[1][]{%
  enhanced,
  colback=bfwBgTeal,
  colframe=bfwPrimaryDark,
  boxrule=1pt, arc=3pt,
  left=10pt, right=10pt, top=6pt, bottom=6pt,
  #1
}

% Glossar-Eintrag
\newcommand{\glossareintrag}[2]{%
  \par\noindent\textbf{\color{bfwPrimaryDark}#1} \enspace #2\par\smallskip
}

% ---------------------------------------------------------------------------
% Abschnitts-Design: farbiges Nummern-Quadrat + feine Linie unter \section,
% dezent farbige \subsection/\subsubsection
% ---------------------------------------------------------------------------
\usepackage{titlesec}
\newcommand{\bfwSecBadge}[1]{%
  \begin{tikzpicture}[baseline={([yshift=-0.62em]n.center)}]
    \node[fill=bfwPrimaryDark, text=white, font=\normalsize\bfseries,
          minimum width=1.55em, minimum height=1.55em, inner sep=1pt,
          rounded corners=1.5pt] (n) {#1};
  \end{tikzpicture}}
\titleformat{\section}
  {\Large\bfseries\color{bfwPrimaryDark}}
  {\bfwSecBadge{\thesection}}{0.6em}{}
  [\vspace{-0.2em}{\color{bfwPrimary!60}\titlerule[0.8pt]}]
\titleformat{\subsection}
  {\large\bfseries\color{bfwPrimaryDark}}
  {\textcolor{bfwPrimary}{\thesubsection}}{0.5em}{}
\titleformat{\subsubsection}
  {\normalsize\bfseries\color{bfwInk}}
  {\textcolor{bfwPrimary}{\thesubsubsection}}{0.4em}{}
\titlespacing*{\section}{0pt}{1.6em}{1em}
\titlespacing*{\subsection}{0pt}{1.2em}{0.5em}

% TikZ-Stil "skizzenhaft" — etwas weicher als Rough.js, aber stabil
\tikzset{
  roughbox/.style = {
    draw=bfwPrimaryDark, thick, fill=bfwBgTeal,
    rounded corners=4pt, inner sep=6pt
  },
  roughaccent/.style = {
    draw=bfwAccent, thick, fill=bfwBgAmber,
    rounded corners=4pt, inner sep=6pt
  },
  roughline/.style = {
    draw=bfwInkSoft, thick, -{Stealth[length=2.5mm]}
  }
}

% Bullet-Symbole (bewusst kein FontAwesome — vermeidet Font-Installations-Probleme)
\newcommand{\faLightbulb}{$\bullet$}
\newcommand{\faBookmark}{$\bullet$}

% ---------------------------------------------------------------------------
% Kopf-/Fußzeile
%   Kopf links:  Wortlogo „Wirtschaftsfachwirt"
%   Kopf rechts: Dokument-Kurztext, gesetzt via \kopfzeile{...}
%   Fuß links:   © Can Siebert · 2026 — Fuß rechts: Seite X von Y
%   Titelseite (Seite 1) bleibt ohne Kopf-/Fußzeile (\handouttitel setzt
%   \thispagestyle{empty}).
% ---------------------------------------------------------------------------
\usepackage{fancyhdr}
\newcommand{\bfwKopfzeileText}{}
\newcommand{\kopfzeile}[1]{\renewcommand{\bfwKopfzeileText}{#1}}
\pagestyle{fancy}
\fancyhf{}
\fancyhead[L]{\small\bfseries\color{bfwPrimaryDark}Wirtschaftsfachwirt}
\fancyhead[R]{\small\color{bfwInkSoft}\bfwKopfzeileText}
\renewcommand{\headrulewidth}{0.6pt}
\renewcommand{\headrule}{{\color{bfwPrimary}\hrule width\headwidth height 0.6pt}}
\fancyfoot[L]{\small\color{bfwInkSoft}\copyright\ Can Siebert\,\textperiodcentered\,2026}
\fancyfoot[R]{\small\color{bfwInkSoft}Seite \thepage\ von \pageref*{LastPage}}
% Auch \thispagestyle{plain} (z.B. durch \maketitle) soll leer bleiben:
\fancypagestyle{plain}{\fancyhf{}\renewcommand{\headrulewidth}{0pt}\renewcommand{\headrule}{}}

% ---------------------------------------------------------------------------
% \handouttitel{Obertitel}{Untertitel}{Kurs-Label}{Termin-Zeile}
%   Einheitlicher professioneller Titelblock: Dachzeile, farbige Headline,
%   Untertitel, farbige Trennlinie und dezente Meta-Zeile (Kurs/Termin/Dozent).
%   Ersetzt \title/\maketitle. Direkt nach \begin{document} aufrufen.
% ---------------------------------------------------------------------------
\newcommand{\handouttitel}[4]{%
  \thispagestyle{empty}%
  \begingroup
  \setlength{\parindent}{0pt}%
  {\bfwlabelfont\color{bfwPrimary}WIRTSCHAFTSFACHWIRT (IHK)\par}
  \vspace{0.55em}
  {\huge\bfseries\color{bfwPrimaryDark}#1\par}
  \vspace{0.5em}
  {\large\color{bfwInkSoft}#2\par}
  \vspace{0.9em}
  {\color{bfwPrimary}\rule{\linewidth}{2pt}}\par
  \vspace{0.55em}
  {\small\color{bfwInkSoft}#3\ \,\textperiodcentered\,\ #4\ \,\textperiodcentered\,\ Dozent: Can Siebert\par}
  \endgroup
  \vspace{1.4em}%
}


\kopfzeile{VWL Lektion 1 \textperiodcentered{} Grundbegriffe \& Preisbildung}

\begin{document}
\handouttitel{Lektion~1: Volkswirtschaftliche Grundbegriffe \& Preisbildung am Markt}%
  {Vom Bedürfnis zum Marktgleichgewicht}%
  {1.1 Volkswirtschaftliche Grundlagen}%
  {Mi, 19.08.2026, 18:00--21:15 Uhr}

\begin{merksatz}[title={\bfseries Worum geht es in dieser Lektion?}]
Die Volkswirtschaftslehre beginnt bei einem Grundkonflikt: Die menschlichen Bedürfnisse
sind (nahezu) unbegrenzt, die Güter zu ihrer Befriedigung aber knapp. Aus dieser Knappheit
folgt die Notwendigkeit des Wirtschaftens. In dieser Lektion lernen Sie die Grundbegriffe
kennen (Wirtschaftssubjekte, ökonomisches Prinzip, Produktionsfaktoren), unterscheiden
Marktarten und Marktformen und verstehen, wie auf einem Markt aus Angebot und Nachfrage
ein Preis entsteht --- das Marktgleichgewicht bei vollständiger Konkurrenz sowie die
Preisbildung bei unvollständiger Konkurrenz. Das ist die Grundlage für den gesamten
Themenblock 1.1 des Rahmenplans und regelmäßig Gegenstand der IHK-Prüfung.
\end{merksatz}

\tableofcontents
\vspace{1em}
\hrule
\vspace{1em}

\section{Warum wir wirtschaften: Knappheit und ökonomisches Prinzip}

Der Ausgangspunkt allen Wirtschaftens ist der Konflikt zwischen \textbf{Güterknappheit}
einerseits und der Tendenz \textbf{unbegrenzter Bedürfnisse} andererseits. Weil nicht alle
Wünsche mit den vorhandenen Mitteln erfüllt werden können, müssen Menschen, Unternehmen
und der Staat \emph{wählen} und \emph{haushalten} --- sie wirtschaften.

\begin{itemize}
  \item \textbf{Bedürfnis:} ein empfundener Mangel mit dem Wunsch, ihn zu beseitigen
    (z.\,B. Hunger, Sicherheit, Anerkennung).
  \item \textbf{Bedarf:} der Teil der Bedürfnisse, der mit \emph{Kaufkraft} ausgestattet ist.
  \item \textbf{Nachfrage:} der Bedarf, der am Markt tatsächlich \emph{wirksam} wird.
\end{itemize}

Wirtschaften heißt, nach dem \textbf{ökonomischen Prinzip} (Rationalprinzip) zu handeln.
Es tritt in zwei Ausprägungen auf:

\begin{itemize}
  \item \textbf{Maximalprinzip:} Mit \emph{gegebenen Mitteln} einen \emph{möglichst großen}
    Erfolg erzielen (z.\,B. mit festem Werbebudget maximale Reichweite).
  \item \textbf{Minimalprinzip:} Ein \emph{gegebenes Ziel} mit \emph{möglichst geringem}
    Mitteleinsatz erreichen (z.\,B. vorgegebene Stückzahl mit minimalen Kosten).
\end{itemize}

\begin{eselsbruecke}
\textbf{Ma}ximalprinzip: \textbf{M}ittel fest, \textbf{a}lles rausholen.
\textbf{Mi}nimalprinzip: Ziel fest, \textbf{mi}t wenig hinkommen.
Beides gleichzeitig („mit wenig Mitteln viel erreichen``) ist \emph{kein} ökonomisches
Prinzip --- ein beliebter Prüfungsfehler!
\end{eselsbruecke}

\subsection{Wirtschaftssubjekte und Wirtschaftssektoren}

Ein \textbf{Wirtschaftssubjekt} ist das kleinste wirtschaftende Grundelement, das
selbstständig wirtschaftliche Entscheidungen treffen kann --- z.\,B. ein Privathaushalt,
ein Unternehmen oder eine staatliche Körperschaft. Für gesamtwirtschaftliche Betrachtungen
werden die Wirtschaftssubjekte zu \emph{Wirtschaftssektoren} zusammengefasst: private
Haushalte, Unternehmen, Staat, (Banken) und das Ausland.

\subsection{Produktionsfaktoren}

Der Produktionsprozess (Leistungserstellung) erfolgt volkswirtschaftlich durch den Einsatz
der \textbf{Produktionsfaktoren}:

\begin{itemize}
  \item \textbf{Arbeit} --- jede auf Erwerb gerichtete körperliche und geistige Tätigkeit,
  \item \textbf{Boden} --- Anbau, Abbau (Rohstoffe) und Standort,
  \item \textbf{Kapital} --- produzierte Produktionsmittel (Maschinen, Anlagen, Gebäude),
  \item zunehmend auch \textbf{Wissen (Know-how)} als eigenständiger vierter Faktor.
\end{itemize}

\begin{merksatz}
Unterscheiden Sie die \emph{volkswirtschaftlichen} Produktionsfaktoren (Arbeit, Boden,
Kapital, Wissen) von den \emph{betriebswirtschaftlichen} (Elementarfaktoren und dispositiver
Faktor) --- Letztere behandeln wir im BWL-Teil des Kurses.
\end{merksatz}

\section{Der Markt: Arten und Formen}

\subsection{Was ist ein Markt?}

Der \textbf{Markt} ist der (reale oder virtuelle) Ort, an dem Angebot und Nachfrage
zusammentreffen und sich der \textbf{Preis} bildet. Seine Aufgaben: Verteilung der Güter,
bestmögliche Versorgung der Marktteilnehmer, Preisbildung. Ziel der Nachfrager ist das
\emph{Nutzenmaximum}, Ziel der Anbieter die \emph{Gewinnmaximierung} (mindestens
Kostendeckung). Für jede Gesellschaft stellen sich drei Grundfragen: \emph{Was} soll
\emph{wie} und \emph{für wen} produziert werden?

\subsection{Marktarten (nach dem Marktgegenstand)}

\begin{center}
\begin{tabular}{p{3.6cm} p{4.2cm} p{6.2cm}}
\textbf{Marktart} & \textbf{Marktgegenstand} & \textbf{Beispiele}\\
\hline
Faktormarkt & Produktionsfaktoren & Arbeitsmarkt, Immobilienmarkt, Sachkapitalmarkt\\
Gütermarkt & Güter & Konsumgüter, Investitionsgüter, Dienstleistungen\\
Geld- und Kapitalmarkt (Finanzmarkt) & Geld, Kredite, Kapitalanlagen & Devisenmarkt, Aktienmarkt\\
\end{tabular}
\end{center}

\subsection{Vollkommener und unvollkommener Markt}

Der \textbf{vollkommene Markt} ist ein \emph{Idealtyp}, der in der Realität nicht
anzutreffen ist (die Börse kommt ihm sehr nah). Bedingungen:

\begin{enumerate}
  \item \textbf{Homogenität der Güter} --- sachlich gleichwertige Güter, keine Unterschiede
    in Qualität, Aufmachung, Verpackung;
  \item \textbf{Fehlen von Präferenzen} --- keine persönlichen, räumlichen oder zeitlichen
    Vorlieben;
  \item \textbf{vollkommene Markttransparenz} --- alle Teilnehmer haben vollständigen
    Überblick über Preise und Konditionen;
  \item \textbf{unendlich schnelle Reaktionen} der Marktteilnehmer (Punktmarkt);
  \item \textbf{polypolistische Marktstruktur} --- viele Anbieter, viele Nachfrager, freier
    Marktzutritt.
\end{enumerate}

Ist \emph{mindestens eine} Bedingung nicht erfüllt, liegt ein \textbf{unvollkommener
Markt} vor.

\subsection{Marktformen (nach der Zahl der Marktteilnehmer)}

\begin{center}
\begin{tabular}{p{2.6cm}|p{3.6cm}|p{3.6cm}|p{3.9cm}}
\textbf{Nachfrager $\downarrow$ / Anbieter $\rightarrow$} & \textbf{viele} & \textbf{wenige} & \textbf{einer}\\
\hline
\textbf{viele} & (zweiseitiges) Polypol \newline \small z.\,B. Wochenmärkte & Angebotsoligopol \newline \small z.\,B. Mobilfunkanbieter & Angebotsmonopol \newline \small z.\,B. regionaler Wasserversorger\\
\hline
\textbf{wenige} & Nachfrageoligopol \newline \small z.\,B. Großmolkereien ggü. Milchbauern & (zweiseitiges) Oligopol \newline \small Spezialmaschinen & beschränktes Angebotsmonopol\\
\hline
\textbf{einer} & Nachfragemonopol \newline \small Staat bei Straßenbau & beschränktes Nachfragemonopol \newline \small Bundeswehr ggü. wenigen Anbietern & (zweiseitiges) Monopol\\
\end{tabular}
\end{center}

\begin{beispiel}
Benzinmarkt: \emph{wenige} Mineralölkonzerne stehen \emph{vielen} Autofahrern gegenüber
$\Rightarrow$ \textbf{Angebotsoligopol}. Apothekenmarkt: \emph{viele} Apotheken,
\emph{viele} Kunden $\Rightarrow$ \textbf{Polypol} (Monopole sind bei Apotheken
\emph{nicht} typisch --- beliebte Quiz-Falle).
\end{beispiel}

\section{Die Nachfrage}

\subsection{Bestimmungsfaktoren der Nachfrage}

Die Nachfrager (Haushalte, aber auch Unternehmen, Staat, Ausland) wollen mit begrenzten
Mitteln ihren Nutzen maximieren. Wichtige Bestimmungsfaktoren der Güternachfrage:

\begin{itemize}
  \item \textbf{Preis des Gutes:} Bei hohem Preis wird i.\,d.\,R. weniger nachgefragt.
  \item \textbf{Preise anderer Güter:} \emph{komplementäre} Güter (ergänzen sich, z.\,B.
    Drucker und Druckerpatronen) und \emph{substitutive} Güter (ersetzen sich, z.\,B.
    Butter und Margarine).
  \item \textbf{Verfügbares Einkommen} (Konsumsumme).
  \item \textbf{Bedürfnisstruktur} (Mode, Demografie, Technologie, Werte).
  \item \textbf{Zukunftserwartungen} bezüglich Preisen und Einkommen.
\end{itemize}

\subsection{Elastizitäten: Wie stark reagieren die Kunden?}
\label{sec:elastizitaet}

Die \textbf{Preiselastizität der Nachfrage} beantwortet eine ganz praktische Frage:
\emph{Wie stark reagieren die Kunden, wenn der Preis sich ändert?} Erhöht ein Bäcker
den Brötchenpreis um 10\,\%, kaufen die Kunden dann fast unverändert weiter (träge
Reaktion) --- oder laufen sie in Scharen zur Konkurrenz (heftige Reaktion)? Genau
dieses Reaktionsausmaß fasst die Elastizität in einer einzigen Zahl zusammen.

\medskip
\noindent
\begin{minipage}[c]{0.55\textwidth}
\begin{formelbox}
\[
\varepsilon \;=\; \frac{\text{Mengenänderung in \%}}{\text{Preisänderung in \%}}
\]
\end{formelbox}
\end{minipage}\hfill
\begin{minipage}[c]{0.41\textwidth}
\small
\textbf{In Worten:} Oben steht die \emph{Reaktion} der Kunden (um wie viel Prozent
ändert sich die nachgefragte \textbf{Menge}?), unten der \emph{Auslöser} (um wie viel
Prozent hat sich der \textbf{Preis} geändert?). Beides immer in Prozent --- nie in
Euro oder Stück!
\end{minipage}
\medskip

Da Preis und Menge normalerweise gegenläufig reagieren (Preis rauf, Menge runter),
ist $\varepsilon$ meist \textbf{negativ}. Für die Interpretation zählt allein der
\textbf{Betrag} $|\varepsilon|$, also die Zahl ohne Vorzeichen:

\begin{center}
\small
\begin{tabular}{p{2.7cm} p{4.0cm} p{3.3cm} p{4.0cm}}
\toprule
\rowcolor{bfwTabHead}
\textbf{Ergebnis} & \textbf{Bedeutung in einfachen Worten} & \textbf{Typische Güter} & \textbf{Preiserhöhung wirkt auf den Umsatz \ldots}\\
\midrule
$|\varepsilon| > 1$ \newline \textbf{elastisch} &
  Kunden reagieren \emph{stark}: Die Menge ändert sich prozentual stärker als der Preis. &
  Luxusgüter, Urlaubsreisen, Markenartikel --- alles, worauf man verzichten oder ausweichen kann &
  \ldots{} \textbf{senkend}: Der Mengenrückgang frisst die Preiserhöhung mehr als auf.\\
\addlinespace
$|\varepsilon| < 1$ \newline \textbf{unelastisch} &
  Kunden reagieren \emph{schwach}: Die Menge ändert sich prozentual weniger als der Preis. &
  Lebensnotwendiges: Grundnahrungsmittel, Medikamente, Strom, Benzin &
  \ldots{} \textbf{steigernd}: Fast alle kaufen weiter --- zum höheren Preis.\\
\addlinespace
$|\varepsilon| = 0$ \newline \textbf{vollkommen unelastisch} &
  Kunden reagieren \emph{gar nicht}: Die Menge bleibt exakt gleich. &
  lebenswichtige Güter ohne jeden Ersatz, z.\,B. Insulin für Diabetiker &
  \ldots{} \textbf{voll steigernd}: Der Umsatz wächst im selben Maß wie der Preis.\\
\bottomrule
\end{tabular}
\end{center}

\begin{beispiel}[title={SCHRITT FÜR SCHRITT: PREISELASTIZITÄT BERECHNEN}]
Ein Medikament kostet 2,00~€; der Preis steigt auf 2,40~€. Die verkaufte Menge sinkt
von 1.000 auf 980 Packungen.
\begin{enumerate}
  \item \textbf{Preisänderung in \% berechnen:} Änderung $= 2{,}40 - 2{,}00 = 0{,}40$~€.
    Bezogen auf den \emph{Ausgangspreis}: $0{,}40 \div 2{,}00 = \mathbf{+20\,\%}$.
  \item \textbf{Mengenänderung in \% berechnen:} Änderung $= 980 - 1.000 = -20$ Packungen.
    Bezogen auf die \emph{Ausgangsmenge}: $-20 \div 1.000 = \mathbf{-2\,\%}$.
  \item \textbf{In die Formel einsetzen:}
    $\varepsilon = -2\,\% \div (+20\,\%) = -0{,}1$, also $|\varepsilon| = 0{,}1$.
  \item \textbf{Interpretieren:} $|\varepsilon| = 0{,}1 < 1$ $\Rightarrow$
    \textbf{unelastisch} --- die Käufer können auf das lebensnotwendige Gut kaum
    verzichten.
  \item \textbf{Umsatz-Probe:} vorher $2{,}00~\text{€} \times 1.000 = 2.000$~€,
    nachher $2{,}40~\text{€} \times 980 = 2.352$~€. Der Umsatz \emph{steigt} --- genau
    wie es die Tabelle für unelastische Güter vorhersagt.
\end{enumerate}
Zum Vergleich: Bei einem Modeartikel mit Preis $+10\,\%$ und Menge $-25\,\%$ wäre
$|\varepsilon| = 2{,}5$ --- \emph{elastisch}, der Umsatz würde sinken.
\end{beispiel}

\begin{achtung}[title={TYPISCHE FEHLER BEI DER PREISELASTIZITÄT}]
\begin{itemize}
  \item \textbf{Vorzeichen falsch gedeutet:} Ein negatives $\varepsilon$ ist bei der
    Preiselastizität \emph{normal} (Preis rauf, Menge runter). Für
    elastisch/unelastisch zählt nur der Betrag $|\varepsilon|$ --- nicht das
    Minuszeichen.
  \item \textbf{Absolute statt prozentuale Änderungen:} „$0{,}40$~€ geteilt durch
    20~Stück`` ist falsch --- beide Änderungen müssen \emph{in Prozent}, jeweils
    bezogen auf den \emph{Ausgangswert}, in die Formel.
  \item \textbf{Zähler und Nenner vertauscht:} Die \emph{Menge} gehört nach oben,
    der \emph{Preis} nach unten. Merkhilfe: erst die Reaktion, dann der Auslöser.
  \item \textbf{Verwechslung mit der Kreuzpreiselastizität:} Die direkte
    Preiselastizität betrachtet \emph{ein} Gut (eigener Preis $\rightarrow$ eigene
    Menge). Sobald \emph{zwei} Güter im Spiel sind (Butter/Margarine), ist die
    Kreuzpreiselastizität gefragt.
\end{itemize}
\end{achtung}

\subsubsection{Kreuzpreiselastizität und Einkommenselastizität --- kurz erklärt}

Die \textbf{Kreuzpreiselastizität} fragt: Wie reagiert die Menge von Gut~A, wenn sich
der \emph{Preis eines anderen Gutes~B} ändert?
\[
\varepsilon_K \;=\; \frac{\text{Mengenänderung bei Gut A in \%}}{\text{Preisänderung bei Gut B in \%}}
\]
Hier trägt das \textbf{Vorzeichen} die Information: \textbf{positiv} $\Rightarrow$
\emph{Substitute} (Butter wird teurer $\rightarrow$ mehr Margarine); \textbf{negativ}
$\Rightarrow$ \emph{Komplemente} (Drucker wird teurer $\rightarrow$ weniger Patronen).

\begin{beispiel}[title={BEISPIEL: KREUZPREISELASTIZITÄT}]
Der Butterpreis (Gut~B) steigt um $10\,\%$; die nachgefragte Menge an Margarine
(Gut~A) steigt um $5\,\%$. Einsetzen:
$\varepsilon_K = +5\,\% \div (+10\,\%) = +0{,}5$. Das \emph{positive} Vorzeichen
zeigt: Butter und Margarine sind \textbf{Substitute}.
\end{beispiel}

Die \textbf{Einkommenselastizität} fragt: Wie reagiert die nachgefragte Menge, wenn
sich das \emph{Einkommen} der Haushalte ändert? Bei normalen Gütern ist sie positiv
(mehr Einkommen $\rightarrow$ mehr Nachfrage), bei Luxusgütern besonders hoch, bei
inferioren Gütern (z.\,B. einfachste Lebensmittel) negativ --- wer mehr verdient,
kauft davon weniger.

\subsection{Nachfragekurve: Bewegung und Verschiebung}

Die Nachfragekurve verläuft \textbf{fallend}: Steigt der Preis, sinkt die nachgefragte
Menge (Bewegung \emph{auf} der Kurve). Ändern sich \emph{andere} Einflussgrößen als der
Preis, \textbf{verschiebt} sich die gesamte Kurve:

\begin{itemize}
  \item \textbf{Linksverschiebung} (Nachfragesenkung): z.\,B. Gut kommt aus der Mode,
    Realeinkommen sinkt, Preis des Substituts fällt, Preis des Komplements steigt.
  \item \textbf{Rechtsverschiebung} (Nachfrageerhöhung): z.\,B. höhere Realeinkommen,
    optimistische Zukunftserwartungen.
\end{itemize}

\section{Das Angebot}

Bestimmungsfaktoren des Güterangebots:

\begin{itemize}
  \item \textbf{Preis des Gutes:} steigender Marktpreis $\Rightarrow$ Anreiz, die
    Produktionsmenge auszuweiten;
  \item \textbf{Kosten} der Produktionsfaktoren (Material-, Lohnkosten);
  \item \textbf{Stand der Technik} (Produktivität).
\end{itemize}

Die Angebotskurve verläuft \textbf{steigend}. Steigende Faktorkosten oder sinkende
Produktivität verschieben sie nach \emph{links} (Angebotssenkung), Kostensenkungen oder
technischer Fortschritt nach \emph{rechts} (Angebotserhöhung).

\section{Preisbildung im Marktgleichgewicht (vollständige Konkurrenz)}

Liegt eine Kombination aus \emph{vollkommenem Markt} und \emph{Polypol} vor, spricht man
vom Modell der \textbf{vollständigen Konkurrenz}. Kein einzelner Teilnehmer kann den Preis
beeinflussen --- alle sind \textbf{Preisnehmer} und \textbf{Mengenanpasser}.

\subsection{Der Preismechanismus}

\begin{itemize}
  \item \textbf{Angebotsüberschuss} (Nachfragedefizit): Bei hohem Preis übersteigt das
    Angebot die Nachfrage $\Rightarrow$ \emph{Käufermarkt}, Anbieter senken die Preise.
  \item \textbf{Nachfrageüberschuss} (Angebotsdefizit): Bei niedrigem Preis übersteigt die
    Nachfrage das Angebot $\Rightarrow$ \emph{Verkäufermarkt}, Preise steigen.
  \item Der Markt tendiert zum \textbf{Gleichgewicht}: Die angebotene Menge entspricht der
    nachgefragten Menge. Die \emph{Gleichgewichtsmenge} ($x^*$) wird zum
    \emph{Gleichgewichtspreis} ($p^*$) abgesetzt --- grafisch der \textbf{Schnittpunkt} von
    Angebots- und Nachfragekurve. Nur dort findet der mengenmäßig größte Umsatz statt.
\end{itemize}

Dieser Anpassungsprozess über Preis und Menge heißt \textbf{Preismechanismus} oder
\textbf{Marktmechanismus}.

\subsection{Konsumenten- und Produzentenrente}

\begin{itemize}
  \item \textbf{Konsumentenrente:} Nutzengewinn der Nachfrager, die bereit gewesen wären,
    \emph{mehr} als den Gleichgewichtspreis zu zahlen.
  \item \textbf{Produzentenrente:} Mehrerlös der Anbieter, die bereit gewesen wären, ihre
    Güter \emph{billiger} als zum Gleichgewichtspreis zu verkaufen.
\end{itemize}

\subsection{Funktionen des Marktpreises}

\begin{itemize}
  \item \textbf{Signalfunktion} (Informationsfunktion): Der Preis spiegelt die Knappheit
    eines Gutes wider.
  \item \textbf{Ausgleichsfunktion}: Preise gleichen Angebots- und Nachfragemenge
    tendenziell aus.
  \item \textbf{Allokationsfunktion} (Lenkungsfunktion): Preise lenken die
    Produktionsfaktoren in die Produktion marktfähiger Güter.
  \item \textbf{Selektionsfunktion} (Sanktionsfunktion): Anbieter mit zu hohen Kosten und
    Nachfrager, die den Preis nicht zahlen können oder wollen, scheiden aus dem Markt aus.
\end{itemize}

\begin{eselsbruecke}
\textbf{SAAS} --- \textbf{S}ignal, \textbf{A}usgleich, \textbf{A}llokation,
\textbf{S}elektion: die vier Funktionen des Marktpreises.
\end{eselsbruecke}

\section{Preisbildung bei unvollständiger Konkurrenz}

In der Realität existieren auf unvollkommenen Märkten unterschiedliche Preise, weil Güter
heterogen sind, Präferenzen bestehen und die Markttransparenz fehlt.

\subsection{Unvollkommenes Polypol: der monopolistische Absatzbereich}

Viele Anbieter stehen vielen Nachfragern auf einem \emph{unvollkommenen} Markt gegenüber.
Durch \textbf{Produktdifferenzierung} (Marken, Werbung, Qualität) schafft sich der Anbieter
einen \textbf{monopolistischen Absatzbereich}: Innerhalb dieses Preisspielraums kann er den
Preis nahezu frei festlegen, ohne größere Nachfrageeinbußen. Überschreitet er ihn, wandern
Kunden ab; unterschreitet er ihn, gewinnt er Kunden, deckt aber evtl. seine Kosten nicht
mehr.

\subsection{Oligopol}

Wenige Anbieter mit großen Marktanteilen (z.\,B. Waschmittel, Automobile, Mineralöl)
beobachten sich gegenseitig. Typische Verhaltensweisen:

\begin{itemize}
  \item \textbf{Marktverdrängungspolitik:} aggressive Preissenkungen, um Marktanteile zu
    gewinnen (ruinöser Preiskampf --- selten);
  \item \textbf{Parallelverhalten:} abgestimmtes Verhalten ohne Absprache, häufig
    \emph{Preisruhe} (Preisstarre) und \emph{Gleichschritt bei Preisänderungen};
  \item \textbf{Preisführerschaft:} ein Anbieter erhöht die Preise, die übrigen ziehen nach
    (z.\,B. Tankstellen vor Feiertagen). Der Wettbewerb verlagert sich auf Qualität, Marke,
    Werbung und Service.
  \item Gefahr verbotener \textbf{Absprachen} (Kartellbildung) --- dazu mehr in Lektion~2
    (Wettbewerbspolitik).
\end{itemize}

\subsection{Monopol}

Beim \textbf{Angebotsmonopol} beliefert nur ein Anbieter den Markt --- er ist theoretisch
\textbf{Preisfixierer}. Die Marktmacht ist besonders stark, wenn kaum Substitutionsgüter
existieren und das Produkt dringend benötigt wird (z.\,B. regionaler Wasserversorger).
Langfristig bestehen Monopole nur bei Marktzutrittsbeschränkungen.

\section{Formeln \& Rechenbeispiele}

Alle Berechnungen dieser Lektion folgen einfachen Grundformeln. Rechnen Sie die
Beispiele einmal selbst nach --- genau solche kleinen Rechnungen kommen in der
IHK-Prüfung vor.

\subsection{Preiselastizität, Kreuzpreiselastizität \& Einkommenselastizität}

Formeln, Interpretations-Tabelle und die durchgerechneten Schritt-für-Schritt-Beispiele
(Medikament: Preis $2{,}00 \rightarrow 2{,}40$~€, Menge $1.000 \rightarrow 980$
Packungen; Butter/Margarine) finden Sie gebündelt in
Abschnitt~\ref{sec:elastizitaet} --- inklusive der Box „Typische Fehler``.
Prüfen Sie sich selbst: Können Sie das Medikamenten-Beispiel ohne Blick in den
Lösungsweg nachrechnen und das Ergebnis interpretieren?

\subsection{Gleichgewichtspreis und -menge aus einer Tabelle bestimmen}

\begin{merksatz}[title={\bfseries Vorgehen}]
Gesucht ist der Preis, bei dem \textbf{angebotene Menge = nachgefragte Menge}. Diese
Zeile liefert Gleichgewichtspreis $p^*$ und Gleichgewichtsmenge $x^*$; der Umsatz im
Gleichgewicht beträgt $p^* \cdot x^*$.
\end{merksatz}

\begin{beispiel}
Für einen Markt liegen folgende Zahlen vor:

\begin{center}
\begin{tabular}{cccl}
\textbf{Preis} & \textbf{Nachfrage (Stück)} & \textbf{Angebot (Stück)} & \textbf{Marktlage}\\
\hline
2~€ & 180 & 60  & Nachfrageüberschuss $\Rightarrow$ Verkäufermarkt\\
3~€ & 140 & 100 & Nachfrageüberschuss $\Rightarrow$ Verkäufermarkt\\
\textbf{4~€} & \textbf{120} & \textbf{120} & \textbf{Gleichgewicht}\\
5~€ & 90  & 150 & Angebotsüberschuss $\Rightarrow$ Käufermarkt\\
\end{tabular}
\end{center}

\begin{enumerate}
  \item \textbf{Zeile suchen}, in der Angebot und Nachfrage übereinstimmen: bei
    \textbf{4~€} sind es jeweils 120 Stück.
  \item \textbf{Ablesen:} Gleichgewichtspreis $p^* = 4$~€, Gleichgewichtsmenge
    $x^* = 120$ Stück.
  \item \textbf{Umsatz im Gleichgewicht:} $4~\text{€} \times 120 = 480$~€.
\end{enumerate}
\textbf{Interpretation:} Oberhalb von $p^*$ (5~€) übersteigt das Angebot die Nachfrage
--- Käufermarkt, Preisdruck nach unten. Unterhalb von $p^*$ (2~€ oder 3~€) übersteigt
die Nachfrage das Angebot --- Verkäufermarkt, die Preise steigen. Der Preismechanismus
führt den Markt zum Gleichgewicht zurück.
\end{beispiel}

\begin{merksatz}[title={\bfseries Wichtig für die Prüfungsvorbereitung}]
Sie sollten sicher können: die Definition „Markt``, die \emph{fünf Bedingungen des
vollkommenen Marktes}, die \emph{Marktformen-Matrix} mit Beispielen, die Skizze der
\emph{Preisbildung} (Gleichgewicht, Überschüsse), die \emph{vier Preisfunktionen},
den Unterschied zwischen \emph{Bewegung auf} und \emph{Verschiebung von} Kurven sowie
die \emph{Rechenwege} für Preiselastizität, Kreuzpreiselastizität und Gleichgewichtspreis
(Abschnitte „Elastizitäten`` und „Formeln \& Rechenbeispiele``).
Üben Sie mit dem Aufgabenheft und dem Quiz dieser Lektion!
\end{merksatz}

\vspace{1em}
\hrule
\vspace{0.8em}
\noindent\textsf{\small Quellen: Rahmenplan Wirtschaftsbezogene Qualifikationen (DIHK), Abschnitt 1.1.1;
Textband Volkswirtschaft (DIHK-Bildungs-GmbH), Kap.~1, S.~1--11. Nächste Lektion (Sa, 22.08.2026):
Wettbewerbspolitik, Eingriffe des Staates in die Preisbildung \& Volkswirtschaftliche Gesamtrechnung.}

\end{document}
